% !TEX root = ../main.tex
\chapter[Supplementary Optimisation Studies]{Supplementary Analyses and Evaluation Protocol}
\label{app:optimisation-supplement}

This appendix preserves the exploratory depth-prediction study and the detailed
evaluation protocol supporting Chapter~5. The first reports completed probe
measurements; the second describes the proposed feedback comparison. The main
chapter presents the completed prompt-search, allocation, and ownership-check
results.

\section{Exploratory Hidden-State Prediction of Assigned Depth}
\label{sec:ch5-depth-probe}

Chapter~3 studies representations of compositional hop count. A separate
probe examines their association with the assigned depth target $d_a$
defined in Section~\ref{sec:ch5-allocation}. The panel contains $2\,648$ reference arms from $1\,000$
MuSiQue questions with BM25 rankings. Supporting passages are assigned to
the earliest-ranking arm; $31$ arm targets belong to questions whose
required evidence is unreachable within the retained lists. The binary
targets are $d_a\leq k$ for $k\in\{2,4,8,16,32\}$.

The frozen models encode the reference subquestion templates as raw text,
without a chat template. Right-padding-aware last-token pooling supplies
the features at nine sampled layers per model. Of the $2\,648$ templates,
$1\,517$ contain unresolved \texttt{\#N} references and differ from the
gold-resolved queries used for retrieval. The representations therefore
describe the templates rather than the executed queries.

\begin{table}[H]
\centering
\caption{Exploratory depth prediction on the fixed reference panel.
Mean AUROC averages the five depth thresholds. Each model's layer is
selected using the same cross-validation results reported here.}
\label{tab:ch5-depth-probe}
\begin{tabular}{lrr}
\hline
Features & Selected layer & Mean AUROC \\
\hline
Qwen2.5-7B-Instruct & 14 & 0.8397 \\
Llama-3.1-8B-Instruct & 12 & 0.8454 \\
Mistral-7B-Instruct-v0.3 & 16 & 0.8442 \\
Surface features and top-ten BM25 scores & n/a & 0.7639 \\
\hline
\end{tabular}
\end{table}

Five-fold cross-validation keeps all arms of a question together. Feature
standardisation is fitted within each training fold. Logistic regression
uses inverse regularisation strength $C=0.01$ for hidden states and $C=1$
for the surface-and-score comparator, so Table~\ref{tab:ch5-depth-probe}
compares two specified readouts as well as their features. Near-constant
hidden-state dimensions are filtered using the full panel before splitting.

The split retains substantial source overlap: $1\,557$ test-arm instances
have an identical template in their training fold, and $755$ test questions
share a source subproblem with training. Layer selection on the reported
folds adds selection optimism. All selected-layer activations are finite;
the Qwen final-layer arrays contain nonfinite activations for $2\,237$ arms,
which the probe code replaces with zero. That layer is unsuitable for
interpreting a layer-wise performance decline.

These measurements motivate a component-separated evaluation of template
representations with nested layer selection and matched readout tuning.
Their use in an allocator requires a further end-to-end comparison of
evidence coverage, answer quality and representation-extraction cost.

\section{Independent Evaluation Design}
\label{sec:ch5-prospective}

\subsection{Frozen Data Roles and Corpus Scope}

The proposed feedback comparison separates questions used for reflection, questions used
for candidate selection and questions reserved for final confirmation.
The source is the official answerable MuSiQue development
release~\cite{trivedi2022musique}. Excluding all composite question IDs in the
historical thousand-question panel leaves $1\,417$ cases. Removing three
additional exact normalised question-text matches leaves $1\,414$.

Of those questions, $363$ have every gold supporting passage in the fixed
$11\,656$-passage corpus. The initial matched-backbone study uses this
eligible subset. Eligibility is determined before new model outputs, by exact
passage identity rather than observed retrieval success. It defines a
corpus-covered target population. Broader MuSiQue performance requires a
separate evaluation with the additional corpus material included.

Questions sharing any underlying single-hop source ID are grouped before
the new splits are assigned. The final manifests contain $96$ optimisation,
$48$ selection and $200$ confirmation questions, with zero shared source
IDs across these three splits. Table~\ref{tab:ch5-splits} gives their
composition. Split identities and file hashes are frozen before candidate
generation.

\begin{table}[H]
\centering
\small
\caption{Frozen question-new splits for the prospective comparison.
Components connect questions sharing a source single-hop ID.}
\label{tab:ch5-splits}
\begin{tabular}{lrrrrr}
\hline
Role & Questions & Two-hop & Three-hop & Four-hop & Components \\
\hline
Optimisation & $96$ & $23$ & $40$ & $33$ & $29$ \\
Selection & $48$ & $16$ & $19$ & $13$ & $22$ \\
Confirmation & $200$ & $73$ & $77$ & $50$ & $65$ \\
\hline
\end{tabular}
\end{table}

All selected composite question IDs and texts differ from the historical
panel and the ownership classifier's training panel. Historical subproblem
reuse remains: $343$ of the $344$ selected questions share at least one
single-hop source with the historical thousand questions. These are new
compositions under the documented exposure exclusions, rather than wholly
new atomic facts. There is no source-ID overlap with the classifier's
training questions. The confirmation set therefore tests transfer to new
compositions within the corpus-covered regime.

The reflection service accepts only hash-bound optimisation samples. The
selection set supplies scores for ranking candidates and remains adaptive
development data. Confirmation question contents and labels stay outside
the search interface until the final instructions and analysis choices have
been fixed.

\subsection{Comparators and Search Allowances}

The fresh study includes the unchanged initial instruction, nonadaptive
instruction sampling, OPRO-style answer-feedback search and three GEPA
feedback conditions. Nonadaptive proposals are generated from the task and
initial instruction before their evaluation results are available. The
OPRO control uses score histories and answer errors. The three GEPA
conditions share the optimizer, reward and execution interface while varying
the feedback view.

Each fresh search is assigned a ceiling of $512$ example-level rollout
evaluations and $24$ candidate proposals. A rollout is an executed
candidate--question pair. Parent comparisons, proposed-candidate comparisons
and validation executions count toward that ceiling. The different methods
allocate their allowance differently, so search progress is measured against
cumulative evaluations and actual cost.

The nonadaptive schedule screens $24$ instructions on a common set of
$12$ optimisation questions, then evaluates its four leaders on the $48$
selection questions. This uses
\begin{equation}
  24\times12+4\times48=480
  \label{eq:ch5-budget}
\end{equation}
candidate--question evaluations. The OPRO control uses four rounds of six
proposals, with a declared twelve-question screen per round and selection
evaluation for each round's nominee. The reflective search uses its allowance
for parent and child batches and complete selection evaluations. The number
of completed mutations is reported explicitly rather than inferred from
the evaluation ceiling.

Within a search condition, the winner is selected by mean selection F1.
Ties are resolved by lower mean execution tokens and then instruction hash.
The fixed initial instruction and the historical OPRO winner are retained
as identified reference policies. The fresh pipeline may use a different
API model configuration from the historical Llama pilot; comparisons are
made within a frozen fresh backbone, and historical scores remain in their
original regime.

The principal contrast is structured versus answer-only feedback. Structured
versus raw feedback is required to isolate the effect of diagnostic
organisation from additional supervision. A second predeclared search seed
provides a replication of search variability when resources permit. Both
seeds' trajectories must then be retained.

\subsection{Outcome Analysis and Cost Accounting}

After selection, each final policy is evaluated on the same locked
confirmation questions. The primary outcome is the question-weighted mean
paired F1 difference. Exact match, complete evidence in the reader contexts,
union depth, masked depth, arm count and model tokens are secondary
measurements. All questions, including fallback executions, contribute to
the aggregate results.

Paired bootstrap resampling preserves each question's outcomes across
policies. A complementary component-level bootstrap keeps questions sharing
a source component together. The confirmation panel contains $65$ such
components, with at most $20$ questions in one component. Question and
component resampling describe sensitivity under different sampling units;
neither supplies a distribution-free generalisation guarantee.

Search cost and inference cost are recorded separately. Search cost includes
instruction proposals, reflection and candidate evaluation. Inference cost
includes decomposition, retrieval-associated model calls and every reader
step for the selected policy. Cached deterministic evaluations count toward
the logical search allowance while their actual charges are recorded
separately. A cache hit is reuse of a result, not an independent execution.

The passage-allocation replay in Chapter~4 provides a complementary cost
diagnostic. On recorded lists it compares balanced total-slot allocation with
the smallest oracle prefix allocation covering the gold passages. The oracle
has access to gold identities and frozen rankings. It measures available
allocation headroom; it is not a learned execution policy. In the sequential
pipeline, changing an early reader allowance may change its answer and every
downstream query. Such a policy requires fresh execution before claiming an
answer-quality or runtime gain.

\subsection{Separating Prompt and Allocation Effects}

The completed studies intervene at different points: the OPRO pilot changes
the decomposition instruction, while Section~\ref{sec:ch5-allocation}
changes prefixes selected from fixed rankings. A prospective factorial
comparison can test their combination. Freeze an initial prompt $p_0$, a
search-selected prompt $p_1$, balanced allocation $g_0$ and learned
allocation $g_1$ before final evaluation. Evaluate all four pairs
$(p_i,g_j)$ on the same questions and reader budget. If $Y_{ij}$ denotes
mean final-answer F1, the interaction is
\begin{equation}
  I=(Y_{11}-Y_{10})-(Y_{01}-Y_{00}).
  \label{eq:ch5-factorial}
\end{equation}
It asks whether the allocation benefit changes under the new prompt.
Each prompt generates its own traces under the same five-passages-per-hop
policy; both allocators then act on those traces. This comparison concerns
pooled context, with trace-generation cost reported separately. Its
outcomes remain prospective.

A sequential extension must choose each arm's budget using information
available at that step and reserve resources for later arms. It then
recomputes intermediate answers and downstream retrieval. The offline
allocator uses resolved later queries and their score profiles, which
would give an online policy future information. Step-local allocation
and re-execution are therefore part of that extension.
